\section{Sola Scriptura in the Middle Ages}

The doctrine of Sola Scriptura was once understood in a total sense: that is, as a call to assimilate one's self to the Logos of God through meditation (often on the Psalms), \& to preach this Gospel to every aspect of Life. The faith was not a matter of private devotion, merely. Ruskin's comments are well worth reading -- he is speaking of Art and of Jerome's method of interpretation. We note that the Protestant theory finds little approval in his pages:

\begin{quotationx}
Much more, must the scholar, who would comprehend in any degree approaching to completeness, the influence of the Bible on mankind, be able to read the interpretations of it which rose into the great arts of Europe at their culmination. In every province of Christendom, according to the degree of art-power it possessed, a series of illustrations of the Bible were produced as time went on; beginning with vignetted illustrations of manuscript, advancing into life-size sculpture, and concluding in perfect power of realistic painting. These teachings and preachings of the Church, by means of art, are not only a most important part of the general Apostolic Acts of Christianity; but their study is a necessary part of Biblical scholarship, so that no man can in any large sense understand the Bible itself until he has learned also to read these national commentaries upon it, and been made aware of their collective weight. The Protestant reader, who most imagines himself independent in his thought, and private in his study, of Scripture, is nevertheless usually at the mercy of the nearest preacher who has a pleasant voice and ingenious fancy; receiving from him thankfully, and often reverently, whatever interpretation of texts the agreeable voice or ready wit may recommend: while, in the meantime, he remains entirely ignorant of, and if left to his own will, invariably destroys as injurious, the deeply meditated interpretations of Scripture which, in their matter, have been sanctioned by the consent of all the Christian Church for a thousand years; and in their treatment, have been exalted by the trained skill and inspired imagination of the noblest souls ever enclosed in mortal clay. \textbf{46}. There are few of the fathers of the Christian Church whose commentaries on the Bible, or personal theories of its gospel, have not been, to the constant exultation of the enemies of the Church, fretted and disgraced by angers of controversy, or weakened and distracted by irreconcilable heresy. On the contrary, the scriptural teaching, through their art, of such men as Orcagna, Giotto, Angelico, Luca della Robbia, and Luini, is, literally, free from all earthly taint of momentary passion; its patience, meekness, and quietness are incapable of error through either fear or anger; they are able, without offence, to say all that they wish; they are bound by tradition into a brotherhood which represents unperverted doctrines by unchanging scenes; and they are compelled by the nature of their work to a deliberation and order of method which result in the purest state and frankest use of all intellectual power. \textbf{47}. I may at once, and without need of returning to this question, illustrate the difference in dignity and safety between the mental actions of literature and art, by referring to a passage, otherwise beautifully illustrative of St. Jerome's sweetness and simplicity of character, though quoted, in the place where we find it, with no such favouring intention,---namely, in the pretty letter of Queen Sophie Charlotte, (father's mother of Frederick the Great,) to the Jesuit Vota, given in part by Carlyle in his first volume, ch. iv.

`` `How can St. Jerome, for example, be a key to Scripture?' she insinuates; citing from Jerome this remarkable avowal of his method of composing books;---especially of his method in that book, \emph{Commentary on the Galatians}, where he accuses both Peter and Paul of simulation, and even of hypocrisy. The great St. Augustine has been charging him with this sad fact, (says her Majesty, who gives chapter and verse,) and Jerome answers, `I followed the commentaries of Origen, of five or six different persons, who turned out mostly to be heretics before Jerome had quite done with them, in coming years, `And to confess the honest truth to you,' continues Jerome, `I read all that, and after having crammed my head with a great many things, I sent for my amanuensis, and dictated to him, now my own thoughts, now those of others, without much recollecting the order, nor sometimes the words, nor even the sense'! In another place, (in the book itself further on\footnote{\url{http://www.gutenberg.org/files/24428/24428-h/24428-h.htm\#Footnote\_3-15\_38}}) he says, `I do not myself write; I have an amanuensis, and I dictate to him what comes into my mouth. If I wish to reflect a little, or to say the thing better, or a better thing, he knits his brows, and the whole look of him tells me sufficiently that he cannot endure to wait.' Here is a sacred old gentleman whom it is not safe to depend upon for interpreting the Scriptures,---thinks her Majesty, but does not say so,---leaving Father Vota to his reflections.'' Alas, no, Queen Sophie, neither old St. Jerome's, nor any other human lips nor mind, may be depended upon in that function; but only the Eternal Sophia, the Power of God and the Wisdom of God: yet this you may see of your old interpreter, that he is wholly open, innocent, and true, and that, through such a person, whether forgetful of his author, or hurried by his scribe, it is more than probable you may hear what Heaven knows to be best for you; and extremely improbable you should take the least harm,---\textbf{while by a careful and cunning master in the literary art, reticent of his doubts, and dexterous in his sayings, any number of prejudices or errors might be proposed to you acceptably, or even fastened in you fatally, though all the while you were not the least required to confide in his inspiration. 48}. For indeed, the only confidence, and the only safety which in such matters we can either hold or hope, are in our own desire to be rightly guided, and willingness to follow in simplicity the guidance granted. But all our conceptions and reasonings on the subject of inspiration have been disordered by our habit, first of distinguishing falsely---or at least needlessly---between inspiration of words and of acts; and secondly by our attribution of inspired strength or wisdom to some persons or some writers only, instead of to the whole body of believers, in so far as they are partakers of the Grace of Christ, the Love of God, and the Fellowship of the Holy Ghost. In the degree in which every Christian receives, or refuses, the several gifts expressed by that general benediction, he enters or is cast out from the inheritance of the saints,---in the exact degree in which he denies the Christ, angers the Father, and grieves the Holy Spirit, he becomes uninspired or unholy,---and in the measure in which he trusts Christ, obeys the Father, and consents with the Spirit, he becomes inspired in feeling, act, word, and reception of word, according to the capacities of his nature. He is not gifted with higher ability, nor called into new offices, but enabled to use his granted natural powers, in their appointed place, to the best purpose. A child is inspired as a child, and a maiden as a maiden; the weak, even in their weakness, and the wise, only in their hour. That is the simply determinable \emph{theory} of the inspiration of all true members of the Church; its truth can only be known by proving it in trial: but I believe there is no record of any man's having tried and declared it vain.\footnote{\url{http://www.gutenberg.org/files/24428/24428-h/24428-h.htm\#Footnote\_3-16\_39}}

Beyond this theory of general inspiration, there is that of special call and command, with actual dictation of the deeds to be done or words to be said. I will enter at present into no examination of the evidences of such separating influence; it is not claimed by the Fathers of the Church, either for themselves, or even for the entire body of the Sacred writers, but only ascribed to certain passages dictated at certain times for special needs: and there is no possibility of attaching the idea of infallible truth to any form of human language in which even these exceptional passages have been delivered to us. But this is demonstrably true of the entire volume of them as we have it, and read,---each of us as it may be rendered in his native tongue; that, however mingled with mystery which we are not required to unravel, or difficulties which we should be insolent in desiring to solve, it contains plain teaching for men of every rank of soul and state of life, which so far as they honestly and implicitly obey, they will be happy and innocent to the utmost powers of their nature, and capable of victory over all adversities, whether of temptation or pain. Indeed, the Psalter alone, which practically was the service book of the Church for many ages, contains merely in the first half of it the sum of personal and social wisdom. The 1st, 8th, 12th, 14th, 15th, 19th, 23rd, and 24th psalms, well learned and believed, are enough for all personal guidance; the 48th, 72nd, and 75th, have in them the law and the prophecy of all righteous government; and every real triumph of natural science is anticipated in the 104th.
\end{quotationx}


Ruskin lived when it was still possible to walk down and view a truly Catholic society in something other than monuments.


\flrightit{Posted on 2012-11-30 by Logres}

\begin{center}* * *\end{center}

\begin{footnotesize}
\begin{sffamily}
\texttt{George on 2013-02-02 at 09:38 said:}

``... and in the measure in which he trusts Christ, obeys the Father, and consents with the Spirit, he becomes inspired in feeling, act, word, and reception of word, according to the capacities of his nature.''

This would be unacceptable today, as it immediately undermines belief in equality and what would follow is Democracy -- why should all have equal vote when all have different abilities to discern ``according to the capacities of his nature''?


\hfill

\end{sffamily}
\end{footnotesize}

